\documentclass[11pt,a4paper]{article}
\usepackage{ftm}

\title{Birch–Swinnerton-Dyer\\[4pt]
       \large Caderno de Lacunas Formais — Programa FTM}
\author{Federated Trivalent Manifold Research Programme}
\date{2026}

\begin{document}
\maketitle
\tableofcontents
\bigskip

%% ---------------------------------------------------------------
\section{Enunciado oficial}
%% ---------------------------------------------------------------

\begin{definition}[Conjectura de Birch–Swinnerton-Dyer]
Seja $E/\mathbb Q$ uma curva elíptica com $L$-função
$L(E,s)=\prod_p L_p(E,s)$, continuada inteiramente.
A conjectura afirma:
\[
  \ord_{s=1}L(E,s)=\rank E(\mathbb Q).
\]
O coeficiente líder satisfaz a fórmula refinada
\[
  \lim_{s\to1}\frac{L(E,s)}{(s-1)^r}
  =\frac{\Omega_E\cdot\mathrm{Reg}(E)\cdot\prod_p c_p\cdot|\Sha(E)|}
         {|E(\mathbb Q)_{\mathrm{tor}}|^2},
\]
onde $\Omega_E$ é o período real, $\mathrm{Reg}(E)$ o regulador de
altura, $c_p$ os números de Tamagawa e $\Sha(E)$ o grupo de
Tate–Shafarevich.
\end{definition}

\bigskip

%% ---------------------------------------------------------------
\section{Programas de pesquisa e encaixe FTM}
%% ---------------------------------------------------------------

\begin{programs}

\item \textbf{Gross–Zagier e Kolyvagin.}
Caso $r\le1$ essencialmente provado:
\begin{itemize}[leftmargin=2em]
  \item Gross–Zagier (1986): $L'(E,1)\neq0\Rightarrow r=1$
    via pontos de Heegner e derivadas de $L$.
  \item Kolyvagin (1988): sistemas de Euler fornecem
    $\Sha(E)$ finito quando $r=0$ ou $r=1$.
\end{itemize}
O caso $r\ge2$ permanece aberto: não existe construção aritmética
sistemática de geradores do grupo de Mordell–Weil.

\begin{gap}
Para $r\ge2$, não há método provado para construir pontos racionais
independentes via teoria analítica de $L(E,s)$. \OPEN
\end{gap}

\item \textbf{Teoria de Iwasawa.}
Sobre a torre ciclotômica $\mathbb Q(\mu_{p^\infty})$,
o $p$-ádico $L$-função $L_p(E,s)$ e o grupo de Selmer
$\Sel_{p^\infty}(E/\mathbb Q^\infty)$ satisfazem conjecturalmente
a Main Conjecture de Iwasawa:
\[
  \mathrm{char}_{\Lambda}\Sel_{p^\infty}=(\mathcal L_p(E)),
\]
onde $\Lambda=\mathbb Z_p[[\Gamma]]$ é a álgebra de Iwasawa.
Skinner–Urban provaram um lado sob hipóteses técnicas.

\begin{gap}
A Main Conjecture completa (igualdade de ideais característicos)
em toda generalidade permanece aberta; em particular para
curvas com redução aditiva ou $p$ de mau comportamento. \OPEN
\end{gap}

\item \textbf{Pontos de Heegner.}
Para discriminante $D<0$ com $L(E,1,\chi_D)\neq0$,
a teoria de corpos de classe fornece um ponto $y_K\in E(K)$
cujo traço $y_K+\bar y_K\in E(\mathbb Q)$ é não trivial sse
$L'(E,1)\neq0$.

A extensão para $r\ge2$ via pontos de Heegner em torres de corpos de
classe de grau superior é uma lacuna ativa (programa Darmon).

\item \textbf{Famílias de $L$-funções e distribuição de postos.}
A conjectura de Goldfeld prevê que a metade das curvas numa família
têm $r=0$ e a outra metade $r=1$. Katz–Sarnak identificam
a distribuição monodrômica. A média de postos $\bar r\le 1/2$
foi estabelecida para certas famílias; $\bar r=1/2$ com controle
refinado está aberto.

\end{programs}

%% ---------------------------------------------------------------
\section{Funtor FTM para curvas elípticas}
%% ---------------------------------------------------------------

\begin{definition}[Funtor $\Phi$]
Busca-se um funtor
\[
  \Phi:\mathrm{Ell}_{\mathbb Q}\longrightarrow\mathrm{FTM}_{\mathrm{arith}}
\]
que associa a cada $E/\mathbb Q$ um objeto
$\Phi(E)=(M_E,\Gamma_E,\mathcal F_E,\lambda_E)$
preservando modelos locais em todos os primos e no lugar arquimediano.
\end{definition}

%% ---------------------------------------------------------------
\section{Obrigações formais FTM para BSD}
%% ---------------------------------------------------------------

\begin{obligation}[BSD1 — Fatores locais]\label{ob:BSD1}
Recuperar $L_p(E,s)$ e o condutor $N_E$ funcionalmente a partir de
$\Phi(E)$, incluindo casos de redução boa, multiplicativa e aditiva.
\end{obligation}

\begin{obligation}[BSD2 — Produto global]\label{ob:BSD2}
Provar que o invariante FTM coincide com $L(E,s)$ como função
de $s\in\mathbb C$, incluindo continuação analítica e equação
funcional.
\end{obligation}

\begin{obligation}[BSD3 — Lado algébrico]\label{ob:BSD3}
Recuperar $E(\mathbb Q)$, grupo de torsão, emparelhamento de altura
de Néron–Tate e regulador $\mathrm{Reg}(E)$.
\end{obligation}

\begin{obligation}[BSD4 — Descida e $\Sha$]\label{ob:BSD4}
Representar os grupos de Selmer $\Sel_n(E)$ e controlar $\Sha(E)$
sem presumir sua finitude quando ela faz parte do conteúdo.
\end{obligation}

\begin{obligation}[BSD5 — Equivalência]\label{ob:BSD5}
Provar que a ordem de anulação de $L(E,s)$ em $s=1$ e o posto
de $E(\mathbb Q)$ coincidem no modelo FTM exatamente com os
objetos clássicos.
\end{obligation}

\begin{remark}
A conexão análoga $\Var(L)\leftrightarrow V_\Omega$ é sugestiva
mas não uma redução: $\Var(L)$ é propriedade global analítica,
enquanto $V_\Omega$ é detector local de defeito.
Grafos de redução em primos isolados não contêm a informação global.
\end{remark}

\MATURITY{M0/M1 conceitual. BSD1–BSD5 todas ABERTAS.}

\end{document}
